% Page budget: 3 pages across the two Section 5 source fragments.
\clearpage
\section{Value Interpretation}\label{sec:value}

Section~4 determines which strategies remain.  This section explains why its
two preservation objects have economic value.  It keeps the value
decomposition, one descendant-event bound, and the finite shadow value needed
to interpret the compression problem.  Appendix~C contains the core
sensitivity formulas; extended belief-occupation and interaction results are
in Online Supplements~S1--S3.  Resource burden is not a third productive
channel: it enters through the outer capacity or penalized problem of
Section~4.

\subsection{Operational--generative decomposition}

Use the unified raw process of \cref{thm:raw-to-compressed-projection}.  Its
\emph{passive value} freezes the library and rules out research:
\begin{align}
 P_0(b,L)&=0,\nonumber\\
 P_{n+1}(b,L)&=F_L(b)+\beta
   \sum_{b'}P_B(b,b')P_n(b',L).\label{eq:passive-value}
\end{align}
Let \(U_n(b,L)\) be the corresponding full value with Continue and every
feasible positive-duration research project, and let
\(\Omega_n(b,L):=U_n(b,L)-P_n(b,L)\) be the research-option premium.

\begin{definition}[Insertion-value channels]
\label{def:innovation-components}
For a candidate \(s\), define
\begin{align}
 \mathcal I_n(s\mid b,L)
   &:=U_n(b,L\cup\{s\})-U_n(b,L),\label{eq:total-innovation}\\
 \Delta^{\op}_n(s\mid b,L)
   &:=P_n(b,L\cup\{s\})-P_n(b,L),\label{eq:operational-innovation}\\
 \Delta^{\gen}_n(s\mid b,L)
   &:=\Omega_n(b,L\cup\{s\})-\Omega_n(b,L).
   \label{eq:generative-innovation}
\end{align}
\end{definition}

\begin{theorem}[Operational--generative value decomposition]
\label{thm:value-decomposition}
For every finite horizon, belief, raw library, and candidate,
\begin{equation}\label{eq:value-decomposition}
 \mathcal I_n(s\mid b,L)
 =\Delta^{\op}_n(s\mid b,L)+\Delta^{\gen}_n(s\mid b,L).
\end{equation}
\end{theorem}

This is an accounting identity, not a sign claim.  The operational term asks
what insertion changes when invention is switched off.  The generative term
is the residual change in the option to research, so it can be negative if
insertion changes costs, completion laws, or the relative appeal of projects.
The projection theorem gives \(U_n(b,L)=V_n(b,K_L)\), allowing both channels
to be evaluated on the sufficient frontier--closure state.

\begin{corollary}[Frontier-silent insertion]
\label{cor:frontier-silent-insertion}
If \(F_{L\cup\{s\}}=F_L\), then
\(\Delta^{\op}_n(s\mid b,L)=0\).  If also
\(C_{L\cup\{s\}}=C_L\), then \(\mathcal I_n(s\mid b,L)=0\).
\end{corollary}

\begin{example}[Operationally silent, generatively valuable bridge]
\label{ex:operational-zero-generative-positive}
A zero-profile bridge carries the unique module enabling a duration-one,
zero-cost project whose descendant pays two.  With one belief, suspended
operation, and \(\beta=1/2\), insertion gives
\begin{equation}\label{eq:operational-zero-generative-positive}
 \Delta^{\op}_2=0,
 \qquad \Delta^{\gen}_2=1.
\end{equation}
The strategy is silent on the current frontier and valuable solely as an
input into future invention.
\end{example}

\subsection{Joint descendant-event lower bound}

The generative channel becomes economically sharp when one retained carrier
opens a specific path to a descendant.  Write \(L^+=L\cup\{s\}\).  For a
project \(q\) of duration \(d\le h\) and a distinguished descendant \(g\),
define the terminal-belief/descendant event mass from the unified joint law:
\begin{equation}\label{eq:t6-joint-descendant-mass}
 \eta_{q,g}(b'\mid b,K)
 :=\Pr\!\left(B_d=b',\ O=\operatorname{some}(g)
       \mid B_0=b,K,q\right).
\end{equation}
This is a subprobability mass and permits arbitrary correlation between the
terminal belief and admission.  The operating-timing adjustment is
\begin{equation}\label{eq:t6-operating-adjustment}
 A^{\op}_{q,h}(b,L^+)
 :=\mathbb E^q_{b,K_{L^+}}\!\left[
   G^{\op}(K_{L^+},q,\mathbf B)
   +\beta^dP_{h-d}(B_d,L^+)\right]-P_h(b,L^+).
\end{equation}
It records operating value earned or forgone while research consumes time.

\begin{theorem}[Joint descendant-event generative-option lower bound]
\label{thm:generative-carrier-lower-bound}
Suppose \(F_{L^+}=F_L\); \(q\), of duration \(d\le h\), is feasible at
\(L^+\) and infeasible at \(L\); its cost is \(\kappa\); and the deleted comparator has
\(\Omega_h(b,L)=0\).  Let \(G(b')\ge0\) lower-bound the complete
full-minus-passive continuation gain after admitting \(g\) on every
length-\(d\) path.  If failure leaves the library unchanged, other admissions
are insertion-only, and full continuation retains the passive policy, then
\begin{equation}\label{eq:t6-lower-bound-adjusted}
 \Delta_h^{\gen}(s\mid b,L)
 \ge \max\!\left\{
  -\kappa+A^{\op}_{q,h}(b,L^+)
  +\beta^d\sum_{b'}\eta_{q,g}(b'\mid b,K_{L^+})G(b'),\,0
 \right\}.
\end{equation}
If \(A^{\op}_{q,h}=0\), this reduces to
\begin{equation}\label{eq:t6-lower-bound}
 \Delta_h^{\gen}(s\mid b,L)
 \ge \max\!\left\{
  -\kappa+\beta^d\sum_{b'}
  \eta_{q,g}(b'\mid b,K_{L^+})G(b'),\,0
 \right\}.
\end{equation}
\end{theorem}

Under the event-specific factorization
\[
 \eta_{q,g}(\,\cdot\mid b,K_{L^+})
 =\pi\rho^d\mu_{q,d}(b,\,\cdot),
\]
write
\(\overline G_{q,d}(b):=\sum_{b'}\mu_{q,d}(b,b')G(b')\).

\begin{corollary}[Product specialization and signs]
\label{cor:t6-independent-product}
\label{cor:t6-comparative-statics}
Under that factorization, the terminal term in
\eqref{eq:t6-lower-bound-adjusted} is
\(\beta^d\pi\rho^d\overline G_{q,d}(b)\).  Holding other terms fixed,
the guarantee rises with pointwise joint mass, pointwise gain, and the
operating adjustment and falls with cost.  Total descendant mass alone has no
sign if its terminal-belief distribution also changes.
\end{corollary}

The joint bound does not require independence.  Duration has no
unconditional sign because it can change
discounting, the joint law, belief occupation, operating rewards,
continuation gains, and horizon feasibility.  Appendix~A gives the proof and
exact retained-carrier example; Online Supplement~S3 gives the
detailed delay algebra.

\subsection{Capacity shadow value}

The hard-capacity problem in \cref{eq:capacity-value} yields the exact
forward shadow value
\begin{equation}\label{eq:capacity-shadow-value}
 \Delta_B^\delta V_\theta^\star(b;B)
 :=V_\theta^\star(b;B+\delta)-V_\theta^\star(b;B),
 \qquad \delta>0.
\end{equation}
By \cref{thm:capacity-value}, this finite difference is nonnegative and is
zero whenever the interval crosses no value-changing attainable burden.  It
is the appropriate marginal object because capacity value is a finite step
function, not a differentiable relaxation.

Scarce capabilities can be complementary.  In the exact two-unit witness,
neither first capability is productive alone but the pair unlocks value one,
so the unit shadow sequence is \(0,1\).  Capacity value therefore need not be
concave and marginal capacity values need not diminish.  Conversely,
additive independent gains with equal unit weights give weakly decreasing
grid shadows.  Any normalized capacity arc additionally requires positive
base capacity and positive base value; when either denominator is zero, we
report \eqref{eq:capacity-shadow-value} itself.
